% Capítulo 4, Desenvolvimento e Resultados Parciais

Este capítulo documenta a construção do ambiente de pesquisa e os achados da etapa atual, estruturando-se a partir da infraestrutura de testes até a análise do comportamento do sistema. A definição da arquitetura experimental estabelece a topologia da rede e os mecanismos de geração de carga, provendo a base para os ensaios práticos. De forma complementar, a especificação formal modela as regras de transição para validar rigorosamente as propriedades do sistema e antecipar divergências de estado. O confronto entre a teoria e a prática ocorre na execução do \emph{model checking} e do \emph{baseline} experimental, etapa que expõe a resposta operacional sob cenários controlados de interrupção e degradação de rede. O conjunto de ferramentas desenvolvidas, incluindo o código de orquestração e os \emph{pipelines} de análise estatística, é consolidado na descrição do estado dos artefatos. O texto atinge então a apresentação dos resultados parciais e das limitações metodológicas, articulando os padrões de violação detectados com as lacunas de avaliação identificadas na revisão da literatura.

\section{Arquitetura experimental}
\label{sec-arquitetura}

O ambiente experimental documentado nesta etapa atende a três requisitos. O primeiro é reproduzir os modos de operação do Infinispan relevantes ao trabalho, em particular o modo \texttt{DIST\_SYNC} para experimentos com replicação síncrona e o modo \texttt{DIST\_ASYNC} para experimentos com replicação assíncrona. O segundo é permitir injeção parametrizada de falhas das classes definidas na Seção~\ref{sec-falhas}, sem alteração no código do protótipo. O terceiro é expor os observáveis necessários para detecção de violações e medição de percentis de latência.

O \emph{cluster} é composto por três nós Infinispan, executados em \emph{containers} Linux gerenciados por \emph{Podman} \citep{podman_docs} com o arquivo de orquestração \texttt{cluster/podman-compose.yml} no repositório \texttt{tcc-desenvolvimento}. Cada nó executa o servidor Infinispan na versão estável, com configuração de \emph{cache} distribuído baseada na documentação oficial e \texttt{numOwners=2}, o que mantém duas cópias por chave \citep{infinispan_docs_stable}. A rede entre \emph{containers} é uma \emph{bridge} controlada pelo \emph{Podman}, com latência e perda de pacotes injetadas via \texttt{tc} no \emph{network namespace} de cada nó, conforme o cenário em execução. Métricas são exportadas em formato OpenMetrics \citep{openmetrics_spec} conforme a documentação de métricas \citep{infinispan_metrics}, e podem ser agente Prometheus opcional \citep{prometheus_docs}.

\revini{DES-01}Na versão atual, os \emph{containers} Podman executam no modo \emph{rootful}, sobre a rede \emph{netavark}, em uma máquina virtual Ubuntu Server~24.04 sob Hyper-V. O ambiente WSL2 usado no início não tinha o módulo de \emph{kernel} \texttt{sch\_netem}, necessário às falhas de rede, e foi substituído pela máquina virtual. A descoberta de membros usa JGroups com transporte TCP e lista explícita de nós, e a pilha de detecção de falhas inclui FD\_ALL3 e FD\_SOCK2 (Seção~\ref{sec-falhas}).\revfim{DES-01}

A geração de carga ocorre por um programa Java executado fora do \emph{cluster}, via Hot Rod \citep{infinispan_docs_stable}, que produz operações concorrentes sobre chaves de sessão. As operações implementadas correspondem ao conjunto definido na Seção~\ref{sec-operacoes}. Para cada operação, o programa de carga registra marca de tempo de início e fim com precisão de nanossegundos, código de retorno, \revini{DES-02}servidor de conexão do cliente\revfim{DES-02} e parâmetros de entrada. As checagens de invariantes ocorrem em dois momentos: imediatamente após cada operação, o programa verifica a coerência local (por exemplo, uma validação que retorna sucesso após \emph{logout} confirmado constitui violação registrada de I1); e periodicamente, em intervalos configuráveis, um auditor consulta todas as réplicas e detecta divergências persistentes (violações de I5 e I6 sob a premissa de quiescência). \revini{DES-03}\revini{DES-62}Na versão atual, a verificação de I5 e I6 é feita por um laço de convergência executado depois que a carga cessa (Seção~\ref{sec-estado-artefatos}).\revfim{DES-62}\revfim{DES-03}

\section{Especificação formal inicial}
\label{sec-spec-formal}

Entre as opções de ferramental discutidas na Seção~\ref{sec-fund-formal}, a escolha adotada nesta etapa inicial é TLA+ com o verificador TLC \citep{yu_lamport_tla_modelchecking_1999}, pela maturidade do ferramental, pelo registro de aplicação industrial a protocolos de replicação \revini{DES-37}\citep{newcombe_aws_tla_2014, newcombe_aws_formal_methods_2015}\revfim{DES-37} e pela adequação da semântica de transições e invariantes ao recorte de invariantes I1 a I6 deste trabalho. Ferramentas como Must, \revini{DES-38}que oferece uma interface de concorrência com vários modelos de comunicação e verifica os programas por redução parcial de ordem dinâmica \citep{enea_must_2024},\revfim{DES-38} \revini{DES-57}ficaram, no planejamento, como alternativa a revisitar caso o refinamento das transições de canal previsto para o Trabalho de Conclusão de Curso~II exigisse primitivas que escapassem a TLA+.\revfim{DES-57} A especificação formal inicial está escrita em TLA+ e versionada como \texttt{spec/SessionStore.tla} no repositório \texttt{tcc-desenvolvimento}. O módulo declara o conjunto de identidades, o conjunto de sessões e dois estados lógicos relevantes (\texttt{Valid}, \texttt{Invalid}), modelados como funções totais sobre os conjuntos correspondentes. O canal de replicação assíncrona é modelado como uma multiset de mensagens pendentes, com regras de transição que entregam uma mensagem qualquer da multiset em qualquer momento, o que captura atraso não determinístico e reordenação. Cada réplica é um processo com seu próprio estado local, e as transições representam, respectivamente, as operações O1 a O7 da Seção~\ref{sec-operacoes}. A abstração corresponde a um modelo de replicação otimista, distinto de algoritmos de consenso baseados em líder eleito como Raft \citep{ongaro_raft_2014}, escolha consistente com o objeto de estudo (replicação assíncrona com reconciliação) declarado no Capítulo~\ref{cap-introducao}. \revini{DES-39}O modelo com falhas, descrito adiante, foi escrito no próprio TLA+, e Must e o tradutor Apalache (Seção~\ref{sec-escolha-ferramenta}) seguem sem uso, como alternativas para ampliar a cobertura.\revfim{DES-39}

\revini{DES-04}Na versão atual, a especificação está acompanhada de um módulo de verificação, \texttt{spec/MCSession.tla}, que declara os cenários e as propriedades verificadas. Cada sessão e cada identidade é um registro único por réplica, com uma marca de tempo ($ts$), espelhando o objeto imutável que o programa de carga grava como unidade no \emph{data grid}. As ações correspondem às operações O1 e O3 a O7 da Seção~\ref{sec-operacoes}. Cada escrita local é difundida às demais réplicas por filas de mensagens pendentes, uma por réplica, e qualquer mensagem pendente pode ser entregue a qualquer momento, o que captura atraso e reordenação. Na entrega, a réplica aplica o registro recebido somente se a marca de tempo da mensagem for maior que a local. Essa regra é um dispositivo de modelagem para evitar regressões arbitrárias no canal abstrato, e não o mecanismo de resolução de conflitos do Infinispan, que sob partição segue a política de fusão configurada.\revfim{DES-04}

\revini{DES-40}O modelo completa o estado com a função auxiliar $\textsf{dono} : \mathcal{S} \to \mathcal{U}$, que associa a cada sessão $s$ a identidade $u$ à qual a sessão pertence. A relação $\textsf{sessions}(u)$, empregada nos invariantes da Seção~\ref{sec-invariantes}, é a inversa de $\textsf{dono}$: $s \in \textsf{sessions}(u)$ se e somente se $\textsf{dono}(s) = u$. No modelo, duas operações diferem da definição da Seção~\ref{sec-operacoes}. O efeito do \emph{login} (O1) inclui $\textsf{dono}'(s) = u$ e não altera o contador, cuja reposição é a operação O5. O bloqueio por limite (O4), por decisão de domínio, não invalida as sessões existentes, e só o bloqueio administrativo (O6) tem cascata de invalidação.\revfim{DES-40}

Os invariantes I1 a I6 são declarados como fórmulas sobre o estado global, com I1 a I4 caracterizadas como propriedades de segurança incondicionais e I5 e I6 como propriedades de convergência condicionadas ao predicado de quiescência (canal de mensagens vazio). A verificação por TLC explora o espaço de estados a partir de configurações iniciais com um número pequeno de identidades, sessões e réplicas (tipicamente duas a três de cada). Em conformidade com a decisão metodológica registrada na Seção~\ref{sec-frente-experimental}, a especificação inicial \revini{DES-58}não modelava\revfim{DES-58} falhas explícitas; a execução do verificador e a leitura dos contraexemplos produzidos integram a Seção~\ref{sec-mc-resultados}. \revini{DES-59}O refinamento da especificação com a inclusão de transições de falha e o tratamento explícito dos contraexemplos obtidos estavam planejados para o Trabalho de Conclusão de Curso~II, e a inclusão das falhas, descrita adiante, foi feita nessa etapa.\revfim{DES-59}

\revini{DES-41}Na versão atual, os invariantes também são declarados sobre a função $\textsf{dono}$, e a quiescência é o estado em que todas as filas de mensagens estão vazias. O registrador de cada réplica modela também $\textsf{dono}$ e $\beta.\textsf{marca}$, e as formas de I5 e I6 foram ampliadas para incluí-los. I5 passa a exigir a convergência do estado lógico e do dono de cada sessão:
\[
\textsf{quiescent}(Q) \Rightarrow \forall s, \forall r_1, r_2 \in \mathcal{R}: \, \sigma_{r_1}(s) = \sigma_{r_2}(s) \, \land \, \textsf{dono}_{r_1}(s) = \textsf{dono}_{r_2}(s).
\]
I6 passa a exigir a convergência do contador e da marca de bloqueio:
\[
\begin{aligned}
\textsf{quiescent}(Q) \Rightarrow \forall u, \forall r_1, r_2 \in \mathcal{R}: \, & \beta_{r_1}(u).\textsf{counter} = \beta_{r_2}(u).\textsf{counter} \\
& \land \, \beta_{r_1}(u).\textsf{marca} = \beta_{r_2}(u).\textsf{marca}.
\end{aligned}
\]
As fórmulas de I1 e I2 da Seção~\ref{sec-invariantes} tratam do estado lógico da sessão. No modelo com várias réplicas, a verificação as complementa com versões avaliadas réplica a réplica, descritas na Seção~\ref{sec-mc-resultados}. I3 e I4 estão declarados no modelo TLA+ em formas de domínio mais fracas, que são consequência da tipagem e aparecem na mesma seção. A forma temporal de I4 é verificada réplica a réplica, e a de I3 não foi verificada.\revfim{DES-41}

\revini{DES-05}A especificação evoluiu em etapas. A versão do Trabalho de Conclusão de Curso~I não modelava falhas. Em setembro de 2026, a entrega passou a difundir cada escrita a todas as réplicas, a comparação de marcas de tempo foi corrigida e a cascata de invalidação do bloqueio administrativo foi refeita. No início de outubro, uma revisão do mecanismo de conflito substituiu as variáveis independentes por campo pelo registro único por chave e retirou o desempate entre réplicas, que deixara de ter função. \revini{DES-63}A revisão se apoiou no código-fonte do Infinispan e do programa de carga. Nele, cada sessão e cada identidade são gravadas como um registro único, com um único número de versão, e a mudança visou reproduzir essa estrutura no modelo.\revfim{DES-63} Em seguida, o modelo ganhou as ações de falha: \texttt{Cair} e \texttt{Reviver}, que representam a queda e a recuperação de uma réplica (F1), e \texttt{Particionar} e \texttt{Mesclar}, que representam a partição e a fusão (F2). Na fusão, o lado com mais réplicas (ou, em empate, qualquer dos lados) copia o registro inteiro de uma de suas réplicas para as do outro lado, sem comparação de marcas de tempo. Duas variantes posteriores, uma com perda de mensagens no canal e outra com relógio atualizado pelo máximo na entrega, foram verificadas em cópias separadas da especificação, ainda não incorporadas à versão principal. A variante com perda tem caráter ilustrativo, porque o modelo não tem retransmissão, e nenhuma das duas corrige o problema de I5 e I6 descrito adiante.\revfim{DES-05}

\section{Execução do \emph{model checking} e estado do \emph{baseline} experimental}
\label{sec-plano-execucao}

A execução das semanas 9 e 10 do cronograma do Trabalho de Conclusão de Curso~I compreende duas atividades. A primeira, verificação formal por \emph{model checking}, foi concluída. A segunda, \emph{baseline} experimental sob falhas selecionadas, foi executada para os cenários F1 e F3 em piloto reduzido, e a sua condução é tratada ao final desta seção. \revini{DES-06}Um diagnóstico dos registros brutos, feito em outubro de 2026, mostrou que essa bateria não mediu o ciclo de vida de sessões, como descrito adiante.\revfim{DES-06}

\subsection*{Verificação formal}
\label{sec-mc-resultados}

A verificação formal foi conduzida em cinco rodadas sucessivas do TLC sobre a especificação descrita na Seção~\ref{sec-spec-formal}. O modelo foi instanciado nas três configurações de cardinalidade crescente parametrizadas em T21 da Tabela~\ref{tab-parametros-baseline}: a configuração reduzida \texttt{MC-small} (duas identidades, três sessões, duas réplicas), a configuração intermediária \texttt{MC-medium} (três identidades, três sessões, duas réplicas) e a configuração ampliada \texttt{MC-full} (três identidades, três sessões, três réplicas). \revini{DES-60}\revini{DES-77}As cinco rodadas, de maio a julho de 2026, usaram só essas três configurações; a configuração \texttt{tiny} foi acrescentada depois, nas verificações de setembro de 2026, com o módulo \texttt{MCSession.tla}.\revfim{DES-77}\revfim{DES-60} Os \emph{scripts} de execução e os arquivos de configuração estão versionados em \texttt{spec/run-tlc.sh}, \texttt{spec/MC-small.cfg}, \texttt{spec/MC-medium.cfg} e \texttt{spec/MC-full.cfg} no repositório \texttt{tcc-desenvolvimento}.

\revini{DES-07}Com a inclusão das falhas no modelo, a verificação dos invariantes I1 a I6 passou a usar quatro configurações de cardinalidade, todas com no máximo uma mensagem pendente por fila, relógio limitado a três e no máximo um evento de falha por execução. A configuração \texttt{tiny} tem uma identidade, duas sessões, duas réplicas e limite de bloqueio igual a dois; \texttt{small} tem duas identidades, três sessões e duas réplicas; \texttt{medium} tem três identidades, três sessões e duas réplicas; e \texttt{full} tem três identidades, três sessões e três réplicas. As saídas brutas do TLC estão em \texttt{runs/tlc/} no repositório. A versão anterior deste texto reproduziu uma expectativa registrada antes da execução, e não o resultado medido. Os números a seguir vêm dos registros de saída do verificador.\revfim{DES-07}

\revini{DES-08}Os invariantes I1 e I2 são violados em traços curtos. I1 é violado em traço de 4 estados nas configurações \texttt{tiny}, \texttt{small} e \texttt{medium} e de 5 estados em \texttt{full}; I2, em traço de 4 estados em \texttt{tiny}, \texttt{small} e \texttt{medium} e de 3 estados em \texttt{full}. Nenhum desses traços contém ação de falha. As violações decorrem da forma das fórmulas, avaliadas sobre o estado de todas as réplicas: basta que uma réplica tenha aplicado a invalidação ou o bloqueio e outra ainda não, situação normal enquanto há mensagens em trânsito. As violações de I1 e I2 nessa forma são, portanto, consequência direta da divergência transitória e não informam sobre o protocolo.

Para separar esse efeito, o módulo de verificação acrescentou propriedades avaliadas réplica a réplica. Essas propriedades foram verificadas em um cenário semeado, com uma identidade, duas sessões já válidas, três réplicas, \emph{login} desabilitado e cotas maiores (até três mensagens por fila, relógio até cinco e dois eventos de falha). A versão de I1 por réplica afirma que uma sessão inválida em uma réplica não passa a outro estado nessa réplica no passo seguinte. No cenário semeado, ela não foi violada sem falha nem com queda, de forma exaustiva com cotas reduzidas e até o tempo-limite com as cotas do plano, e foi violada sob partição, em traço de 4 estados. Esse resultado é coerente com a perda, na fusão, de escritas feitas na partição menor entre a divisão e sua detecção, descrita pela documentação do Infinispan \cite[seção 7.1.1]{infinispan_configuring_15}, com duas ressalvas: no modelo, a escrita da réplica isolada é anterior à partição e ainda não propagada, aproximação daquela janela, e a fusão do modelo é uma abstração da política real. O comportamento real sob partição com carga ainda não foi medido.

No modelo com \emph{login} habilitado, a versão de I1 por réplica também é violada sem falha e com queda, por artefatos do modelo. O relógio de cada réplica não acompanha as marcas recebidas, e, na variante que corrige esse ponto, a violação vem de \emph{logins} concorrentes da mesma sessão em réplicas diferentes. Nenhum dos dois ocorre no programa de carga, que usa a versão do registro lido e identificador único de sessão. A versão de I2 por réplica, segundo a qual uma identidade bloqueada não tem sessão válida na mesma réplica, \revini{DES-61}vale para qualquer tipo de bloqueio e é violada em traços curtos, de 3 a 6 estados conforme o cenário. Todos os 28 contraexemplos obtidos vêm do bloqueio por limite (O4), por exemplo no traço formado por um \emph{login} seguido de dois incrementos de tentativas falhas, porque esse bloqueio não invalida as sessões existentes (Seção~\ref{sec-spec-formal}).\revfim{DES-61} Esse caso é tratado como violação esperada por desenho, e a análise experimental o separa do bloqueio administrativo. A versão de I4 por réplica, que exige que o motivo de invalidação não mude, também é violada.

Os invariantes I5 e I6 são violados nas quatro configurações, em traços de 5 a 7 estados, por dois mecanismos. Em cinco dos oito traços mínimos, duas escritas concorrentes na mesma chave recebem a mesma marca de tempo, porque cada réplica atribui marcas pelo próprio contador, e a regra de entrega, que exige marca estritamente maior, faz cada réplica recusar o registro da outra. A divergência persiste depois que as filas se esvaziam. Esse é um defeito de modelagem, surgido com a retirada do desempate entre réplicas, e sua correção, por reintrodução de um desempate ou por canal com entrega em ordem, está pendente. Nos outros três traços (I5 em \texttt{tiny} e I6 em \texttt{small} e em \texttt{full}), a fusão copia o registro do lado vencedor sobre uma escrita local ainda não propagada, e as mensagens dessa escrita, entregues depois, alcançam só as outras réplicas, o que deixa divergência permanente. Esse segundo mecanismo é um resultado do modelo sob partição, coerente com a seção 7.1.1 da documentação citada acima e independente da correção do empate. Até essa correção, os contraexemplos de I5 e I6 que não passam pela fusão não informam sobre o protocolo.

As formas de I3 e I4 declaradas no modelo, contador dentro da faixa de 0 ao limite e motivo dentro do conjunto declarado, não foram violadas. A verificação foi exaustiva apenas na configuração \texttt{tiny}, com 570\,089 estados distintos, 11\,243\,221 estados gerados e profundidade 21. Nas configurações maiores, as execuções foram interrompidas pelo tempo-limite sem violação. Em \texttt{small}, por exemplo, a verificação de I3 percorreu 30\,474\,782 estados distintos em cerca de 9 minutos, e em \texttt{full}, 5\,466\,292 estados distintos em cerca de 4 minutos. Essas formas decorrem da tipagem e não verificam a monotonicidade do contador nem a imutabilidade do motivo. Esta última, avaliada réplica a réplica, é violada, como descrito acima.\revfim{DES-08}

Os traços produzidos pelo TLC alimentam o catálogo inicial de combinações operação $\times$ canal $\times$ invariante que orienta a fase subsequente. O refinamento da especificação para reduzir os contraexemplos por mecanismos de confirmação em duplo envio ou \emph{anti-entropy} dirigida fica registrado como ação para o Trabalho de Conclusão de Curso~II (Apêndice~\ref{ap-plano-tcc2}, etapa 2). \revini{DES-09}Na versão atual, a pendência da especificação é a correção do empate de marcas de tempo descrita acima, e o catálogo passou a relacionar operação, falha e invariante, base do alinhamento entre o modelo e os experimentos previsto no Apêndice~\ref{ap-plano-tcc2}.\revfim{DES-09}

\subsection*{\emph{Baseline} experimental}
\label{sec-baseline-estado}

O \emph{baseline} experimental parcial previsto para o Trabalho de Conclusão de Curso~I cobre dois cenários de falha: \revini{DES-32}falha de nó (F1), executada como parada ordenada,\revfim{DES-32} e atraso com \emph{jitter} no canal (F3). Os parâmetros operacionais de cada cenário (tamanho do \emph{cluster}, modo de replicação, \texttt{numOwners}, distribuição de chaves, mistura de operações, volume e duração) estão consolidados na Tabela~\ref{tab-parametros-baseline} da Seção~\ref{sec-parametros}, com as fontes que os fundamentam. Nesta etapa, os dois cenários foram executados na mesma bateria, contra a mesma condição \emph{sem falha} de referência, e seus resultados são reportados nesta subseção.

A execução do cenário F3 exige o módulo de \emph{kernel} \texttt{sch\_netem} para a injeção de atraso por \texttt{tc-netem} \citep{netem_manpage}, ausente no ambiente WSL2 utilizado no piloto inicial de F1. Para viabilizar os dois cenários sob o mesmo ambiente, a bateria foi reexecutada em uma máquina virtual Linux (Hyper-V, Ubuntu Server~24.04), com o \emph{cluster} Infinispan em contêineres \emph{Podman} no modo \emph{rootful} sobre a rede \emph{netavark}. O atraso de canal de F3 foi aplicado a partir do \emph{host} da máquina virtual, no \emph{network namespace} do contêiner-alvo \texttt{isn2}, acessado via \texttt{nsenter}, com a disciplina \texttt{tc-netem} configurada em \texttt{delay 20ms} e desvio de $13$\,ms sob \texttt{distribution normal}, calibrada para deslocar o percentil 99 em aproximadamente $50$\,ms, o alvo de T18 (sob a aproximação $p99 \approx \text{média} + 2{,}326 \cdot \text{desvio}$). \revini{DES-10}A falha F1 usou \texttt{podman stop} e \texttt{podman start} sobre o mesmo contêiner \texttt{isn2}, isto é, uma parada ordenada, e não a queda silenciosa definida na Seção~\ref{sec-falhas}.\revfim{DES-10} \revini{DES-11}\revfim{DES-11} \revini{DES-42}\revini{DES-75}Com \texttt{SIGTERM}, o nó sai de forma ordenada, mecanismo explicado na Seção~\ref{sec-revisoes}, e por isso essa bateria e as verificações intermediárias que usaram o mesmo comando são descritas como parada ordenada.\revfim{DES-75}\revfim{DES-42}

O \emph{cluster} foi composto por três nós Infinispan, com imagem fixada em \texttt{quay.io/infinispan/server:15.0.21.Final}, \revini{DES-12}em \texttt{DIST\_SYNC} e \texttt{numOwners=2}. A configuração do \emph{cluster} só tinha \emph{caches} síncronos, e todas as execuções na máquina virtual até outubro de 2026, incluindo as seis células desta bateria, ocorreram nesse modo.\revfim{DES-12} A descoberta de membros adotou JGroups com transporte TCP e protocolo \texttt{TCPPING} sobre lista explícita de \emph{hosts}. A formação do \emph{cluster} em três membros foi validada pelo \emph{endpoint} REST de saúde antes de cada repetição, e o nó parado em F1 retornou ao tamanho três entre repetições. A bateria permanece um piloto reduzido em relação à bateria plena desenhada nos parâmetros T9 e T10 (30 repetições de $10^6$ operações). O piloto cobriu seis células, combinação dos dois cenários de mistura S1 ($50/50$) e S2 ($95/5$) com as três condições \emph{sem falha}, F1 e F3, em três repetições de $90$\,s por célula, com janela de injeção de \emph{jitter} de $65$\,s nas células F3. Os eventos do aquecimento foram descartados em todas as células.

\revini{DES-13}Um diagnóstico dos registros brutos, feito em outubro de 2026, mostrou que a bateria não mediu o ciclo de vida de sessões. Em todas as seis células, 100\,\% dos \emph{logins} retornaram recusa por bloqueio, e 100\,\% das validações retornaram ausência de resultado (código \texttt{NONE}), que encobria a leitura de chave ausente; na célula S1 sem falha, por exemplo, os 1\,291\,332 \emph{logins} registrados foram todos recusados. Nenhum evento registrou código de erro. \emph{Login} e validação somaram 75\,\% dos eventos em S1 e 98\,\% em S2. Quatro causas, todas de desenho da carga e da análise, explicam o resultado. A primeira é a identidade única de cada \emph{thread}, que atingia o limite de tentativas e permanecia bloqueada, já que a mistura não continha desbloqueio. \revini{DES-43}Essa identidade servia ao \emph{login} e às operações O4 e O5. Com limiar de bloqueio igual a cinco, O5 preservando a marca de bloqueio e O7 ausente da mistura, o tempo esperado de absorção da cadeia do contador indica bloqueio permanente depois de cerca de 130 operações da \emph{thread} em S1 e 500 em S2, em média, o que corresponde a frações de segundo no ritmo medido. Nessa cadeia, o incremento tem probabilidade 0,10 e a reposição 0,05 em S1; em S2, o incremento tem probabilidade 0,01, sem reposição.
% Fonte: wt-plano-ctrl/docs/diagnostico-login-blocked-2026-10-22.md, secao 2.3 (o documento diz "simulacao", sem script preservado). Valor esperado exato da cadeia, conferido pelo Bibliotecario em 03/11/2026: S1 p_inc = 0,10, p_reset = 0,05 -> 1055/8 = 131,875 operacoes; S2 p_inc = 0,01 sem reset -> 500. Pesos de Scenario.java e limite 5 de SessionOps.java no commit 51a31f0. [DES-74]
\revfim{DES-43}A segunda é a leitura, pelas validações e \emph{logouts} isolados, de chaves sintéticas que nenhuma operação gravava. A terceira é a ausência de limpeza do \emph{cluster} entre execuções, de modo que as identidades já estavam bloqueadas no primeiro evento. A quarta é o mascaramento: o programa registrava a chave ausente com o mesmo código de ausência de resultado, e o \emph{script} de percentis contava recusas e ausências como sucesso, o que ocultou o problema nos resumos.

As latências medidas nessa bateria correspondem, portanto, à recusa de \emph{login}, à leitura de chave ausente e às escritas do contador sobre identidades já bloqueadas, e não às operações que os rótulos indicam. Por isso, a versão anterior deste capítulo, que apresentava os percentis de latência dessa bateria por operação e condição de falha, foi revista: as tabelas foram retiradas, e nenhum número de latência de junho é usado como resultado ou será comparado com os da nova medição. A ausência de erros e de violações nessa bateria também não indica \emph{failover} sem efeito, porque nenhuma sessão foi criada. A carga foi corrigida nos pontos descritos na Seção~\ref{sec-carga}: limpeza e conferência dos \emph{caches} antes de cada execução, identidades de \emph{login} separadas das que recebem tentativas falhas, chaves de sessão gravadas antes da medição, códigos distintos para sessão válida, inválida e ausente, e um verificador de sanidade que recusa execuções incoerentes com a mistura declarada. \revini{DES-44}Os limiares desse verificador, definidos no diagnóstico da bateria de junho, foram fixados no programa de carga congelado antes do piloto da nova medição, de modo que não foram ajustados depois de observados os dados.\revfim{DES-44} A nova medição de referência, com essa carga, foi pré-registrada (Seção~\ref{sec-protocolo-rebaseline}) e ainda não foi executada.\revfim{DES-13}

Os artefatos do piloto estão organizados em \texttt{runs/} com um diretório por célula (\texttt{S1-none}, \texttt{S1-F1}, \texttt{S1-F3}, \texttt{S2-none}, \texttt{S2-F1} e \texttt{S2-F3}), cada qual com o CSV de eventos crus por repetição, o \texttt{summary.json} por célula e o \texttt{comparacao.json} por contraste, gerados pelo \emph{pipeline} de análise em \texttt{analysis/} (\emph{numpy} e \emph{scipy}).

A bateria plena, com os parâmetros T9 e T10 da Tabela~\ref{tab-parametros-baseline} (30 repetições de $10^6$ operações) e a carga em $75\%$ do pico do piloto sem falhas, permanece como protocolo desenhado e escalável a partir do piloto aqui reportado. \revini{DES-14}Esse protocolo deu lugar à nova medição de referência pré-registrada (Seção~\ref{sec-protocolo-rebaseline}), que mede \texttt{DIST\_SYNC} como controle e \texttt{DIST\_ASYNC} como alvo na mesma campanha, em dez repetições de 300\,s por célula, com F1 por queda silenciosa, e que ainda não foi executada.\revfim{DES-14} Nessa bateria, a duração da indisponibilidade do nó em F1 fica fixada em $60$\,s (T17), valor superior ao \emph{timeout} padrão do detector de falhas FD\_ALL3 do Infinispan ($40$\,s, T16), o que garante a detecção do evento pelo \emph{cluster} antes do término da janela. \revini{DES-15}O atraso de canal de F3 segue calibrado para deslocar o percentil 99 em $50$\,ms (T18), via \texttt{tc-netem} \citep{netem_manpage}, valor de engenharia, sem fonte empírica, e calibrável.\revfim{DES-15} Os dois cenários previstos como ampliação opcional, partição (F2) e \emph{state transfer} induzido (F4), \revini{DES-64}ficaram reservados, no planejamento, para o Trabalho de Conclusão de Curso~II.\revfim{DES-64} \revini{DES-16}\revini{DES-65}Esse plano foi modificado, e F4 foi redefinida (Seção~\ref{sec-revisoes}). Os injetores de F2 e F4 passam pelas verificações sem carga de aplicação descritas a seguir, e a campanha com o cenário S3 os exercita sob carga.\revfim{DES-65}\revfim{DES-16}

\subsection*{Verificação dos injetores na máquina virtual}
\label{sec-verificacao-injetores}

\revini{DES-17}Antes de usar os injetores de F2 e F4 sob carga, cada um passa por uma sequência de verificações na máquina virtual, sem carga de aplicação, em que cada etapa só começa depois que a anterior é validada. Para F2, o nó \texttt{isn3} foi isolado dos outros dois por uma janela de cerca de 90\,s em quatro repetições. O tempo entre o fim da janela e o retorno do \emph{cluster} ao estado saudável com três nós foi medido por consultas a cada 2\,s. Nas três repetições feitas em \emph{cluster} recém-criado e com monitoramento mínimo (a primeira, a terceira e a quarta), esse tempo foi de 48\,s, 46\,s e 44\,s. A segunda repetição, feita em um \emph{cluster} que já sofrera duas partições, com uma sonda que escrevia durante a partição e com coleta de monitoramento mais intensa, levou 188\,s. Os registros mostram que o atraso ocorreu depois da fusão da vista, na resolução de conflitos de um dos \emph{caches}, que esgotou o tempo-limite de 60\,s. A repetição seguinte, em \emph{cluster} novo, enfraqueceu a hipótese da sonda, mas as duas hipóteses restantes, histórico do \emph{cluster} e coleta intensa, não foram separadas. A causa não foi isolada, e a repetição é registrada como limitação e reportada à parte. Esses tempos dizem respeito apenas à recomposição da topologia: sem carga de aplicação, nada indicam sobre divergência de dados ou perda de escritas.

Das oito etapas da sequência mais recente, cinco foram executadas. A primeira preparou, no hospedeiro, o pacote de código e seus resumos criptográficos. A segunda, na máquina virtual e sem alterar a rede, instalou esse pacote com conferência dos resumos, verificou as ferramentas e carregou os módulos do \emph{kernel} usados por F4, executou na própria máquina virtual as suítes de teste dos injetores de F2 e F4 e leu o estado do \emph{cluster}, que estava parado depois de uma queda do hospedeiro, com investigação da causa. A terceira verificou que o \emph{kernel} aceita as disciplinas e o filtro de F4 sem afetar o tráfego, com filtro para um endereço reservado a testes. A quarta verificou que o temporizador do sistema armado pelo \emph{script} de verificação dispara e aciona a limpeza do injetor de F4. A quinta repetiu a partição em \emph{cluster} recém-criado. Essa repetição substituiu outra, planejada para o \emph{cluster} com histórico, que perdeu o sentido quando uma queda do hospedeiro desfez o \emph{cluster}, e seu tempo de recomposição está entre os citados acima. Faltam a primeira injeção real de F4, com perda de 20\,\% por 60\,s e limpeza posterior, e duas verificações do injetor de F2, uma da recomposição manual depois de uma partição real e outra do temporizador de segurança, armado pelo \emph{script} de verificação, disparando sozinho. O modo silencioso de F1 foi ensaiado apenas no \emph{cluster} local, e seu ensaio na máquina virtual precede a recriação do \emph{cluster} com os quatro \emph{caches}.\revfim{DES-17}

\section{Estado dos artefatos}
\label{sec-estado-artefatos}

Os artefatos versionados no repositório \texttt{tcc-desenvolvimento} cobrem, ao final do Trabalho de Conclusão de Curso~I, três linhas de trabalho integradas, complementadas por um \emph{pipeline} de análise estatística e por \emph{scripts} de orquestração.

A primeira linha de trabalho é a configuração do \emph{cluster} Infinispan. O arquivo \texttt{cluster/podman-compose.yml} orquestra três nós da imagem oficial fixada em \texttt{quay.io/infinispan/server:15.0.21.Final} sobre uma rede \emph{bridge} dedicada, e o arquivo \texttt{cluster/infinispan-cluster.xml} \revini{DES-18}descreve os \emph{caches} com \texttt{numOwners=2} e 64 segmentos, com parametrização baseada nas linhas T1 a T4, T15 e T16 da Tabela~\ref{tab-parametros-baseline}, exceto pelo número de segmentos, que em T4 é 256. Até outubro de 2026, a configuração tinha apenas os \emph{caches} \texttt{sessions} e \texttt{counters}, ambos síncronos (\texttt{DIST\_SYNC}).\revfim{DES-18} A política de partição adotada é \texttt{DENY\_READ\_WRITES} com \emph{merge-policy} \texttt{PREFERRED\_NON\_NULL}\revini{DES-34} \cite[seção 7]{infinispan_configuring_15}\revfim{DES-34}, que define o comportamento padrão diante de partição involuntária no escopo do Trabalho de Conclusão de Curso~I, dado que a injeção controlada de F2 \revini{DES-66}estava reservada, no planejamento, para o Trabalho de Conclusão de Curso~II. O injetor de F2 foi implementado depois (Seção~\ref{sec-verificacao-injetores}).\revfim{DES-66}

\revini{DES-35}A versão atual acrescenta \texttt{sessions-async} e \texttt{counters-async} em modo assíncrono, com os mesmos parâmetros, selecionados pela opção \texttt{-{}-replication} do programa de carga. Os quatro \emph{caches} são não transacionais. Essa versão foi testada no \emph{cluster} local e só será instalada na máquina virtual quando o \emph{cluster} for recriado para a nova medição de referência.\revfim{DES-35}

A segunda é o programa de carga em Java sobre Hot Rod, organizado em projeto Maven \citep{maven_apache} (Java~17), com sete componentes:

\begin{itemize}
    \item as operações O1 a O7 padronizadas sob o tipo único de retorno \texttt{OpResult};
    \item um gerador de chaves Zipfian autocontido, baseado em amostragem por rejeição-inversão de complexidade $O(1)$ por amostra, com parâmetros $\rho = 0{,}99$ e universo de $100\,000$ chaves, conforme T5, T6 e T8;
    \item um seletor de cenários que codifica as misturas S1 (\emph{50/50}) e S2 (\emph{95/5}) de T7;
    \item uma política de aquecimento (T11) que estende a janela de não-medição em $\max(60\,\text{s},\ 10\,\% \cdot \text{duração})$ seguidos de $5\,\%$ do tempo restante;
    \item instrumentação de latência baseada em \emph{HdrHistogram} \citep{hdrhistogram} com despejo em CSV (T13);
    \item um auditor de invariantes que combina verificação imediata após cada operação para I1 a I4 com auditoria periódica em janelas de quiescência para I5 e I6;
    \item uma interface de linha de comando com onze opções nomeadas que parametriza T9, T10, T11 e T14 e produz mensagem de ajuda autogerada com validação por contrato.
\end{itemize}

\revini{DES-19}Até outubro de 2026, o programa de carga recebeu os seguintes acréscimos:

\begin{itemize}
    \item nas operações O1 a O7, a invalidação do \emph{logout} e da cascata de bloqueio por escrita condicional à versão e códigos que distinguem erro de transporte (\texttt{ERROR\_TRANSPORT}) de recusa por indisponibilidade (\texttt{ERROR\_UNAVAILABLE});
    \item nas misturas S1 e S2, a versão de \emph{slots} e gerações, que é a usada nas medições (Seção~\ref{sec-carga});
    \item o cenário S3, com servidores de aplicação ligados a nós distintos, usuários virtuais, ator administrativo e conjunto de identidades para bloqueio por limite (Seção~\ref{sec-carga-s3});
    \item uma fase de preparação que esvazia e confere os \emph{caches} e grava a carga inicial de \emph{slots};
    \item o registro de cada evento em CSV, com marcas de tempo de início e fim em nanossegundos, código de retorno, servidor de conexão do cliente e, quando há marcadores de falha, fase da execução;
    \item no auditor, a classificação de I2 por marcos de bloqueio e um laço de convergência para I5 e I6, \revini{DES-67}descrito a seguir;\revfim{DES-67}
    \item um verificador de sanidade que lê o CSV de eventos e emite veredito sobre a mistura observada;
    \item o pacote do controlador heurístico, descrito adiante, e opções nomeadas de linha de comando para duração, repetições, aquecimento, modo de replicação, mistura, cenário e semente.
\end{itemize}\revfim{DES-19}

\revini{DES-45}O laço de convergência roda depois que a carga cessa, porque I5 e I6 só são verificáveis sob quiescência. A cada 5\,s, suas rodadas comparam o estado das réplicas, inclusive o dono de cada sessão ($\textsf{dono}_{r_1} = \textsf{dono}_{r_2}$), até a convergência ou até um tempo-limite. O tempo de convergência vai do fim da carga até a primeira de três rodadas consecutivas sem divergência entre réplicas, com teto de 600\,s. A nova medição de referência não usa o laço, e as regras de classificação e de contagem que ela adota estão na Seção~\ref{sec-carga}.\revfim{DES-45}

A terceira é a especificação TLA+ \texttt{spec/SessionStore.tla}, com as operações O1, O3, O4, O5, O6 e O7 e os invariantes I1 a I6 declarados, acompanhada\revini{DES-68}, no Trabalho de Conclusão de Curso~I,\revfim{DES-68} de três configurações de cardinalidade (T21) e de um \emph{runner} de execução. A operação O2, por ser leitura pura, não tem ação dedicada no modelo, conforme registrado na Seção~\ref{sec-spec-formal}. \revini{DES-20}Na versão atual, a especificação inclui as ações de falha e é acompanhada do módulo de verificação \texttt{spec/MCSession.tla} e de quatro configurações de cardinalidade, descritos nas Seções~\ref{sec-spec-formal} e \ref{sec-plano-execucao}.\revfim{DES-20}

Os artefatos complementares incluem o \emph{pipeline} de análise estatística em Python (\texttt{analysis/percentis.py}, \texttt{analysis/bootstrap\_ic.py} e \texttt{analysis/compare\_scenarios.py}, com base em \emph{numpy} \citep{numpy_harris_2020} e \emph{scipy} \citep{scipy_virtanen_2020}, cobrindo T13 e T22) e quatro \emph{scripts} de orquestração: \texttt{scripts/inject-crash.sh} para a falha F1 conforme T17, \texttt{scripts/inject-jitter.sh} para a falha F3 conforme T18 (com remoção idempotente da disciplina \texttt{tc-netem} via \emph{trap} em \texttt{EXIT}, \texttt{INT} ou \texttt{TERM}), \texttt{scripts/collect-metrics.sh} para coleta periódica do \emph{endpoint} OpenMetrics nos três nós conforme T19 e T20 \revini{DES-69}(escrito, mas não integrado, como registrado adiante)\revfim{DES-69}, e \texttt{scripts/run-baseline.sh} como orquestrador da bateria $\text{cenário} \times \text{falha} \times \text{repetição}$. Cada \emph{script} aceita o argumento \texttt{-{}-dry-run} que imprime o plano e os comandos previstos sem dependência de Podman, de \emph{cluster} ativo ou da \emph{capability} \texttt{NET\_ADMIN}.

\revini{DES-21}Na versão atual, os \emph{scripts} de injeção cobrem as quatro classes de falha: \texttt{inject-crash.sh} (F1, nos modos de parada ordenada e de queda silenciosa), \texttt{inject-partition.sh} (F2), \texttt{inject-jitter.sh} (F3) e \texttt{inject-netdegrade.sh} (F4, com disciplina de prioridade, \texttt{netem} por banda e filtro por endereço de destino). As falhas de rede usam \texttt{tc-netem} \citep{netem_manpage}, aplicado a partir do hospedeiro no \emph{network namespace} do \emph{container}, porque a imagem do Infinispan não traz a ferramenta \texttt{tc}. Os injetores de F2 e F4 gravam marcadores de início e fim, removem as disciplinas de rede ao terminar, por rotina de limpeza que roda também quando recebem sinal de término, e oferecem um modo de limpeza avulso. Nas verificações na máquina virtual, os \emph{scripts} de verificação armam ainda um temporizador do sistema que aciona essa limpeza se o injetor for morto. Completam o conjunto o coletor de recursos \texttt{collect-resources.sh} (Seção~\ref{sec-energia}), os orquestradores \texttt{run-baseline.sh} e \texttt{run-rebaseline.sh}, o gerador da matriz da nova medição e os \emph{scripts} de análise em Python (\texttt{rebaseline.py}, \texttt{fases.py}, \texttt{referencia\_p99.py}, \texttt{energia\_proxy.py} e \texttt{comparar\_bracos.py}).\revfim{DES-21} \revini{DES-46}Das quatro classes de falha com injetor, as três primeiras também aparecem no modelo formal, F1 e F2 como ações e F3 como atraso no canal (Seção~\ref{sec-spec-formal}).\revfim{DES-46}

\revini{DES-47}O coletor das métricas secundárias expostas pelo Infinispan (Seção~\ref{sec-metricas}), \texttt{scripts/collect-metrics.sh}, foi escrito em junho de 2026, mas existe apenas em ramos separados do repositório e nunca foi integrado ao ramo principal nem à versão atual, de modo que não foi usado nas baterias. O coletor de recursos mede o uso de CPU de cada \emph{container}, fora do Infinispan. Como a máquina virtual não dá acesso aos contadores de energia do processador, o custo energético é estimado a partir desse tempo de CPU e de um modelo de potência declarado \citep{fan_power_provisioning_2007, kansal_vm_power_2010}, conforme a Seção~\ref{sec-energia}. O coletor de recursos e o \emph{script} de estimativa de energia estão implementados e testados com dados simulados, mas nenhum consumo foi medido em campanha, e a nova medição de referência não coleta uso de CPU.\revfim{DES-47}

A vertente formal foi executada em cinco rodadas do TLC nas três configurações da Tabela~\ref{tab-parametros-baseline} (T21)\revini{DES-70}\revini{DES-78}; a configuração \texttt{tiny} foi acrescentada depois, nas verificações de setembro de 2026. Os resultados estão descritos na Seção~\ref{sec-mc-resultados}.\revfim{DES-78}\revfim{DES-70} A vertente experimental compila e empacota como \emph{jar shaded} de aproximadamente $18\,\text{MB}$ e dispõe de suíte de testes JUnit \citep{junit5}\revini{DES-79}, com seleção determinística\revfim{DES-79} da \emph{seed} para evitar variação amostral entre rodadas. Uma verificação determinística da suíte, conduzida sobre dez rodadas consecutivas em ambiente Linux com OpenJDK Temurin~17 e Maven~3.9, registrou aprovação em todas as rodadas. \revini{DES-22}Na versão atual, o \emph{jar} está congelado para a nova medição, e os componentes têm suítes de testes automatizados em Java, Python e \emph{shell}, estas com simuladores dos comandos de sistema para rodar sem \emph{cluster}.\revfim{DES-22} O \emph{pipeline} de análise estatística foi validado em conjunto sintético de $30\,000$ operações distribuídas entre sete tipos de operação, com detecção, \revini{DES-23}pelo teste de Mann-Whitney por percentil\revfim{DES-23}, do efeito-na-cauda esperado para F1 em relação ao \emph{baseline}. O programa de carga foi executado contra o \emph{cluster} vivo nos cenários F1 e F3, em máquina virtual Linux com o módulo \texttt{sch\_netem} disponível, com o \emph{pipeline} de análise consumindo os CSV crus de cada repetição, conforme reportado na Seção~\ref{sec-baseline-estado}. \revini{DES-24}O diagnóstico dessa execução, descrito na mesma seção, mostrou que ela não mediu o ciclo de vida de sessões.\revfim{DES-24}

\subsection*{Controlador heurístico: projeto e implementação}
\label{sec-controlador}

\revini{DES-48}O controlador heurístico previsto no Trabalho de Conclusão de Curso~I (Seção~\ref{sec-fund-adapt}) foi projetado e implementado em parte na etapa seguinte, e as configurações pré-definidas tomaram a forma de perfis entre os quais ele alterna. Na terminologia MAPE-K, o monitoramento corresponde ao agregador de janelas e à sonda de topologia, a análise à comparação dos sinais com limiares, o planejamento à máquina de estados que escolhe o perfil e a execução à aplicação do perfil aos parâmetros do cliente, componentes descritos a seguir. Os contraexemplos produzidos pelo TLC (Seção~\ref{sec-mc-resultados}) orientam a escrita do programa de carga e o desenho do controlador.\revfim{DES-48}

\revini{DES-25}O controlador executa como \emph{thread} dentro do processo Java do programa de carga, externo ao \emph{data grid}: não roda no servidor, não reconfigura o servidor, apenas lê a saúde dos nós pela API REST, e não grava estado replicado. O núcleo de decisão fica em um pacote que não depende do Infinispan nem do código de sessão, restrição conferida por teste, de modo que o mesmo código pode decidir ao vivo ou sobre registros gravados. O núcleo usa como tempo o fim de cada janela de observação, e não o relógio do sistema, o que o torna determinístico e testável sem espera real. No desenho aprovado, cada janela gerará uma linha em arquivo JSONL, com os observáveis, os limiares em vigor, a decisão e seus motivos, como trilha de auditoria, e o controlador terá três modos, selecionados por opção de linha de comando: desligado, em que nenhum objeto novo é criado; \emph{shadow}, em que observa e decide sem atuar; e \emph{active}, em que aplica o perfil decidido. A trilha e os modos dependem da ligação ao laço de carga, ainda não implementada.

O monitoramento agrega os eventos do cliente em janelas contíguas de 5\,s e consulta a saúde de cada nó a cada 2\,s. O Quadro~\ref{qd-sinais-controlador} resume os sinais. O sinal de latência é a fração de eventos da janela com latência acima do p99 de referência da operação, obtido de execuções sem falha do mesmo cenário e modo de replicação. Para uma única operação, a fração ultrapassa 1\,\% exatamente quando o p99 da janela supera o de referência. Somada sobre operações, ela funciona com poucas amostras por operação. Janelas com amostra abaixo do mínimo tornam o sinal indeterminado. As violações de I5 e I6 ficam fora dos gatilhos, porque só são verificáveis sob quiescência, e entram na avaliação como desfecho.

\begin{table}[htbp]
\captionsetup{type=table, name=Quadro}
\caption{Sinais observados pelo controlador heurístico}
\label{qd-sinais-controlador}
\centering
\footnotesize
\begin{tabular}{p{2.6cm}p{6.4cm}p{4.2cm}}
\hline
\textbf{Sinal} & \textbf{Definição na janela} & \textbf{Papel} \\
\hline
Latência ($S_\text{lat}$) & fração de eventos com latência acima do p99 de referência da operação; mínimo de 100 eventos & primário; exige 2 janelas seguidas \\
Erro ($S_\text{err}$) & fração de eventos com \texttt{ERROR\_TRANSPORT} ou \texttt{ERROR\_UNAVAILABLE}; mínimo de 50 eventos & primário; exige 2 janelas seguidas \\
Topologia ($S_\text{top}$) & consulta de saúde falhou, menos de três nós ou algum \emph{cache} fora do estado saudável & primário; basta 1 amostra \\
Violação ($S_\text{viol}$) & novas violações de I1, I4 ou I2 classificada como real & secundário; severidade máxima \\
\hline
\end{tabular}
\fonte{O autor}
\end{table}

A decisão é uma máquina de dois estados, NORMAL e CONSERVADOR, avaliada uma vez por janela por quatro regras. Pela regra R1, uma violação leva ao perfil conservador imediatamente, mesmo durante o intervalo mínimo entre transições (\emph{cooldown}), porque, em segurança, reagir tarde custa mais que oscilar uma vez. Pela regra R2, fora do \emph{cooldown}, uma anomalia de topologia, ou latência ou erro acima do limiar de entrada em duas janelas seguidas, leva ao perfil conservador. Pela regra R3, no perfil conservador, qualquer sinal alto zera a contagem de janelas calmas. Pela regra R4, o retorno ao perfil normal exige, ao mesmo tempo, \emph{cooldown} expirado, seis janelas calmas seguidas, topologia estável há um intervalo mínimo e nenhuma violação recente. A histerese vem de limiares de saída menores que os de entrada. Os valores atuais são provisórios, fornecidos por arquivo de parâmetros ou pelo conjunto provisório padrão, sem valor fixo na lógica de decisão: limiar de latência de 0,05 na entrada e 0,02 na saída, de erro de 0,02 e 0,005, \emph{cooldown} de 30\,s, estabilidade de topologia de 30\,s e retenção após violação de 120\,s. A calibração final seguirá regra fixada antes de ver os dados de calibração.

O perfil conservador atua por quatro parâmetros do cliente. O primeiro limita a admissão, reduzindo a 50\,\% o número de operações em curso. O segundo reduz o ritmo, com pausa de 5\,ms entre operações em S1 e S2 e tempo de pensar multiplicado por 1,5 em S3. O terceiro adia as escritas administrativas de S3 (bloqueio, desbloqueio e alterações do contador) até a saída do perfil conservador ou até 30\,s, sem descartá-las. O quarto dobra a espera base entre retentativas de S3, com variação aleatória de $\pm 50\,\%$, sem mudar o número de retentativas. Todos apenas escolhem quando o cliente emite uma operação que já emitiria: nenhuma operação é criada ou suprimida, nenhum estado replicado é tocado e a configuração do servidor não muda. Essa restrição mantém os comportamentos do sistema dentro dos admitidos pelo modelo formal e preserva o modo de replicação como variável do experimento. O modo de replicação e o tempo-limite do cliente Hot Rod ficaram fora desta primeira etapa porque alteram o sistema medido ou exigem reconstruir o cliente.

Estão implementados e testados na versão local do código que integra esses componentes, com relógio simulado e capturas reais da saúde do \emph{cluster}: o núcleo de decisão, com as regras, a histerese e o \emph{cooldown}; o agregador de janelas, ligado ao registro de eventos e ao auditor por um ponto de escuta opcional que, desligado, deixa a saída do programa idêntica à anterior; a sonda de topologia; o controle de emissão de operações, com os quatro parâmetros; o coletor de recursos; e os \emph{scripts} de análise por fase, de referência de p99, de estimativa de energia e de comparação pareada entre braços. Falta ligar o controlador ao laço de carga: a interface de linha de comando ainda não tem a opção que o ativa, e por isso o controlador não foi executado com o programa de carga e não há resultado sobre seu efeito. Também faltam a reprodução das decisões sobre registros gravados, os marcadores de início e fim de F1 e F3 nos registros, uma proteção contra parâmetros de energia inválidos, o pré-registro com execução curta de teste, a calibração e a avaliação descrita na Seção~\ref{sec-avaliacao-controlador}.

\label{sec-controlador-cluster}
A atuação do controlador sobre a configuração do \emph{cluster}, prevista no objetivo específico (4), é uma segunda etapa planejada, posterior à ligação do controlador ao laço de carga e à avaliação da atuação no cliente. Ela vem depois por três razões. Durante a nova medição de referência e a campanha com o cenário S3, a configuração do servidor é variável do experimento e não pode mudar. Só a avaliação da etapa cliente indica se atuar no servidor acrescenta algo. Por fim, parte das ações possíveis produz comportamentos que o modelo formal não admite.

A etapa começa por uma verificação, na documentação e em \emph{cluster} local, de quais atributos do Infinispan 15 podem ser alterados em tempo de execução sem reiniciar o \emph{cache}. Os candidatos iniciais, hipótese de partida ainda não verificada, são o tempo-limite de chamadas remotas e o tempo-limite de aquisição de trava, que só alteram quando uma operação termina; a suspensão do rebalanceamento, que altera o momento da convergência; e a mudança forçada da disponibilidade de uma partição, que permitiria escrita na partição minoritária, comportamento fora do modelo. O modo de replicação, o número de cópias, o número de segmentos e a política de partição e de fusão \revini{DES-36}devem exigir a recriação do \emph{cache}, com perda dos dados, o que a mesma verificação vai confirmar, e ficam fora.\revfim{DES-36} A primeira etapa permanece restrita ao cliente e é avaliada como tal; a segunda terá pré-registro e avaliação próprios. A lista final de parâmetros, a admissão de ações que alteram a convergência ou a semântica sob partição, estas com extensão do modelo TLA+, e o tempo de máquina dependem de decisão posterior do autor.\revfim{DES-25}

\section{Resultados parciais e limitações conhecidas}
\label{sec-resultados-parciais}

O resultado parcial central desta etapa é a coerência entre o conjunto de invariantes alvo, o material teórico levantado pelo mapeamento sistemático e os contraexemplos produzidos pela verificação formal. A síntese inicial do mapeamento sustenta a contribuição pretendida do trabalho em duas dimensões complementares. A primeira diz respeito à métrica de avaliação. \revini{DES-26}Nos cerca de trinta trabalhos lidos na íntegra até esta etapa, a taxa de violação de invariantes não aparece como métrica primária de avaliação experimental, pois a literatura se concentra em métricas de desempenho agregadas, como latência média, vazão e percentis de cauda, e em métricas de disponibilidade do sistema; a classificação dos 165 estudos do mapeamento ainda precisa confirmar essa constatação.\revfim{DES-26} A segunda diz respeito ao método. A especificação formal acoplada a invariantes específicos do domínio de sessão não é tratada como objeto central no mesmo conjunto de fontes, e o ponto comparativo mais próximo encontrado, \citet{nguyen_weaker_consistency_2019}, combina consistência relaxada com monitoramento e recuperação reativos, sem verificação formal a montante. Os contraexemplos de I1 e I2 obtidos pelo TLC, comentados na Seção~\ref{sec-mc-resultados}, fornecem o catálogo inicial de padrões de violação que o controlador heurístico previsto para o Trabalho de Conclusão de Curso~II deve mitigar. \revini{DES-27}Na forma global, porém, as violações de I1 e I2 são consequência direta da divergência transitória entre réplicas e não informam sobre o protocolo; as propriedades avaliadas réplica a réplica, descritas na mesma seção, separam esse efeito.\revfim{DES-27} \revini{DES-49}Na etapa atual, além da inclusão de falhas no modelo formal (Seção~\ref{sec-spec-formal}), os injetores das quatro classes de falha foram implementados e a carga foi corrigida depois do diagnóstico da primeira bateria.\revfim{DES-49} \revini{DES-28}O controlador está projetado e implementado em parte, ainda sem ligação ao laço de carga (Seção~\ref{sec-controlador}).\revfim{DES-28} \revini{DES-50}Sua avaliação está desenhada (Seção~\ref{sec-avaliacao-controlador}), mas não foi pré-registrada nem executada. O pré-registro, com os valores finais, vai preceder qualquer execução.\revfim{DES-50}

Três limitações conhecidas devem ser registradas. A primeira é a execução parcial dos artefatos: A vertente formal completou as cinco rodadas planejadas, e o \emph{baseline} experimental teve os cenários de \revini{DES-33}parada ordenada de nó (F1)\revfim{DES-33} e de atraso com \emph{jitter} de canal (F3) executados em piloto reduzido sobre máquina virtual Linux, \revini{DES-29}cujo diagnóstico, na Seção~\ref{sec-baseline-estado}, mostrou que a bateria não mediu o ciclo de vida de sessões\revfim{DES-29}; a bateria plena nos parâmetros T9 e T10 segue como protocolo a escalar a partir do piloto\revini{DES-30}, e a nova medição de referência que a substitui foi pré-registrada e ainda não foi executada\revfim{DES-30}. \revini{DES-51}O pré-registro do protocolo \citep{nosek_preregistration_2018}, feito em outubro de 2026 e acrescido de dois complementos antes de qualquer execução, não admite ajuste depois de observados os dados. O piloto da nova medição, que precede a campanha, também não foi executado.\revfim{DES-51} A segunda é a abstração da rede no modelo formal, que\revini{DES-71}, no Trabalho de Conclusão de Curso~I, capturava atraso e reordenação mas não modelava perdas permanentes nem partições, escopo da evolução prevista para o Trabalho de Conclusão de Curso~II.\revfim{DES-71} \revini{DES-31}Na versão atual, o modelo inclui queda, recuperação, partição e fusão como ações, e a perda de mensagens foi verificada apenas em cópia separada da especificação, sem retransmissão (Seção~\ref{sec-spec-formal}).\revfim{DES-31} A terceira é a execução do mapeamento sistemático em uma única plataforma agregadora, com a complementação por busca manual nas bases originais formalizada entre as ameaças à validade na Seção~\ref{sec-ms} e integrada às atividades remanescentes para o Trabalho de Conclusão de Curso~II.

\subsection*{Revisões em relação ao planejamento}
\label{sec-revisoes}

\revini{DES-52}O plano da Proposta e do Trabalho de Conclusão de Curso~I sofreu revisões durante a execução. Esta subseção registra, para cada revisão, o plano, a mudança e o motivo, e remete às seções com os detalhes técnicos. A renovação de sessão constava do objetivo da Proposta, mas não entrou nas sete operações (O1 a O7, Seção~\ref{sec-operacoes}) definidas no recorte do Trabalho de Conclusão de Curso~I. Como ficou fora do escopo desde esse recorte, a ausência não resulta de revisão durante a execução, e nem o modelo formal nem o programa de carga têm operação de renovação. Os objetivos específicos (Seção~\ref{sec-objetivos-especificos}) foram revistos em três pontos. O objetivo (3) passou a prever a arquitetura experimental nos dois modos de replicação, síncrono e assíncrono, com injeção de queda de nó, partição, atraso e degradação parcial da rede. O objetivo (4) foi dividido em duas etapas: a primeira atua apenas sobre parâmetros do cliente, e a atuação sobre a configuração do \emph{cluster} é uma segunda etapa planejada, posterior à avaliação da primeira (Seção~\ref{sec-controlador-cluster}). O objetivo (5) passou a incluir, além das métricas previstas, o tempo de CPU por operação útil, com a estimativa de energia correspondente, ausente do plano original (Seção~\ref{sec-energia}).\revfim{DES-52}

\revini{DES-53}Na verificação formal, o plano previa um modelo sem falhas no Trabalho de Conclusão de Curso~I e a inclusão de falhas no Trabalho de Conclusão de Curso~II. As duas etapas foram cumpridas. As revisões de estrutura feitas ao longo da execução, com seus motivos, estão na Seção~\ref{sec-spec-formal}, e seus efeitos sobre os resultados de I1 a I6, na Seção~\ref{sec-mc-resultados}.\revfim{DES-53}

\revini{DES-54}Na vertente experimental, quatro revisões alteram a leitura dos dados. A primeira é o modo de replicação: a linha T1 da Tabela~\ref{tab-parametros-baseline} fixa \texttt{DIST\_ASYNC} como alvo, mas todas as execuções na máquina virtual até outubro de 2026 ocorreram em \texttt{DIST\_SYNC} (Seção~\ref{sec-baseline-estado}), achado que levou aos \emph{caches} assíncronos (Seção~\ref{sec-estado-artefatos}) e à medição dos dois modos na mesma campanha. A segunda é a carga, que na bateria de junho não criou sessões e foi corrigida (Seção~\ref{sec-carga}) antes de uma nova medição de referência pré-registrada. O cenário S3, acrescentado em outubro de 2026 e externo à nova medição, exercita o bloqueio administrativo e a escrita em nó isolado (Seção~\ref{sec-carga-s3}). A terceira é a falha F1. Em outubro de 2026, antes de qualquer execução da campanha, a queda silenciosa descrita na Seção~\ref{sec-protocolo-rebaseline} substituiu o \texttt{podman stop} usado na bateria de junho e nas verificações intermediárias, \revini{DES-72}porque a parada ordenada não exercita a detecção de falhas. Com \texttt{SIGTERM}, o servidor fecha o canal JGroups e pede a saída ao coordenador, que instala a nova visão sem que nenhum detector suspeite do nó \cite[seções 3.10.13 e 6.8.7]{jgroups_manual_5}. Na queda silenciosa, a interface de rede é desativada antes do \texttt{SIGKILL} porque, com as conexões TCP fechadas pelo sistema operacional, o FD\_SOCK2 detectaria a queda em cerca de um segundo, e a detecção passaria a não depender do FD\_ALL3 (Seção~\ref{sec-falhas}).\revfim{DES-72} O estado do ensaio desse modo está na Seção~\ref{sec-verificacao-injetores}. A quarta é o conjunto de falhas, tratado adiante.

O protocolo da nova medição também se afasta da Tabela~\ref{tab-parametros-baseline} em pontos declarados. \revini{DES-76}A justificativa de $10^6$ operações por execução e de 30 repetições, dada na tabela, refere-se ao desenho original e não se aplica ao protocolo pré-registrado em outubro de 2026.\revfim{DES-76} As dez repetições por célula substituem as 30 de T10, por limite de tempo de máquina, com a execução como unidade de análise. Essa redução diminui o poder estatístico dos contrastes entre células, de modo que diferenças pequenas entre condições podem não ser detectadas. A duração fixa de 300\,s substitui o volume de $10^6$ operações de T9 e as durações de T14, e o aquecimento mínimo passa a 30\,s (T11). A bateria de junho e as verificações intermediárias usaram \texttt{numSegments} igual a 64, e não 256 como em T4, e a nova medição mantém esse valor. Ele veio da configuração inicial do \emph{cluster}, sem fonte que o justifique, e foi mantido para preservar a comparabilidade entre as baterias. A carga não segue o alvo de 75\,\% do pico de T12: o programa de carga opera em laço fechado, com quatro \emph{threads} que emitem a operação seguinte assim que a anterior termina, sem limitação de vazão. A vazão de cada execução é medida e reportada, e não imposta. Os afastamentos de T4 e de T12 foram declarados em complemento ao registro do protocolo, antes da coleta. Na nova medição, a análise de T22 passa a reamostrar execuções no \emph{bootstrap} e a aplicar o teste de Mann-Whitney, como análise secundária, sobre os p99 por execução, e não sobre eventos agrupados de várias execuções.\revfim{DES-54}

\revini{DES-55}A quarta revisão da vertente experimental é o conjunto de falhas. As três classes de evento planejadas na Introdução (Capítulo~\ref{cap-introducao}) deram lugar a quatro classes: além da queda de um nó seguida de \emph{failover} e da partição de rede, a terceira é o atraso com \emph{jitter} no canal entre os nós, e a quarta é a degradação parcial da rede, definida adiante, no lugar da mudança de topologia com \emph{state transfer} do modelo de falhas planejado (Seção~\ref{sec-falhas}). F2 deixou de ser reservada e ganhou injetor. Em F3, a distribuição lognormal do plano (T18) deu lugar à normal na execução, porque o \texttt{tc-netem} traz tabelas prontas para distribuições como a normal e a de Pareto, mas não para a lognormal, que exigiria gerar e instalar um arquivo próprio no ambiente. A calibração está descrita na Seção~\ref{sec-baseline-estado}, e a lognormal fica como análise de sensibilidade. Essa calibração é do atraso por pacote, e não da latência das operações, que ainda não foi medida sob F3 com carga válida.

Em outubro de 2026, o autor decidiu testar, além da partição total, a perda de pacotes e a degradação parcial, e o rótulo F4 passou a designar a degradação parcial da rede. \revini{DES-73}Falhas de rede reais nem sempre cortam toda a comunicação: ao lado de partições completas, o levantamento de \citeonline[p.~50--53]{bailis_kingsbury_network_2014} relata perdas de 50\,\% a 100\,\% dos pacotes, perturbações que afetam só alguns grupos de nós e uma placa de rede que descartava os pacotes recebidos, mas não os enviados.\revfim{DES-73} Na nova definição de F4, a comunicação entre o nó alvo e seus pares sofre perda, atraso, duplicação ou reordenação de pacotes, de forma simétrica ou assimétrica, sem corte completo. O tráfego do cliente não é afetado, porque o filtro casa apenas os endereços dos pares do \emph{cluster}. As severidades de perda vão de N1 a N5 (1\,\%, 5\,\%, 20\,\%, 50\,\% e 99,9\,\%), e a mais alta deve aproximar uma partição, hipótese a confirmar experimentalmente. O rebalanceamento e o \emph{state transfer}, que alteram o mapeamento entre chaves e nós \citep{infinispan_docs_stable}, continuam presentes como consequência da reentrada do nó em F1 e da fusão em F2, mas não têm injetor próprio.\revfim{DES-55}

\revini{DES-56}No controlador, o plano previa um processo externo que lesse observáveis em arquivo e passasse ao perfil conservador quando a taxa de violação superasse um limiar. O projeto aprovado em outubro de 2026 o executa como \emph{thread} dentro do processo do programa de carga (Seção~\ref{sec-controlador}), porque o observável primário é por evento e a serialização para outro processo custaria CPU, que é métrica avaliada. O gatilho primário passou a combinar latência de cauda relativa, taxa de erro e topologia, e as violações ficaram como gatilho secundário, por três motivos. Primeiro, I5 e I6 só são verificáveis sob quiescência. Segundo, o bloqueio por limite produziria registros de I2 esperados por desenho (Seção~\ref{sec-mc-resultados}). Terceiro, nas paradas de nó da bateria de junho nenhum erro foi registrado, o que sugere que a detecção dependa de latência e topologia. Essa observação tem valor apenas indicativo, porque a bateria não criou sessões (Seção~\ref{sec-baseline-estado}). A atuação sobre a configuração do \emph{cluster}, prevista na Proposta e no objetivo (4), ficou para a segunda etapa mencionada acima.

Os critérios de avaliação do controlador também mudaram. O desenho aprovado em outubro de 2026 mantém a margem de 10\,\% e a rastreabilidade, mas troca a unidade de análise, do evento para a execução, e substitui o teste de Mann-Whitney, previsto na Seção~\ref{sec-metricas}, pelo desenho em blocos pareados da Seção~\ref{sec-avaliacao-controlador}. O teste aplicado a eventos trata como independentes observações de uma mesma execução e é sensível a diferenças no corpo da distribuição, não só na cauda. O teste de Mann-Whitney sobre os p99 por execução permanece apenas na nova medição de referência, como análise secundária.\revfim{DES-56}
